\documentclass[11pt,a4paper]{article}

% ── packages ──────────────────────────────────────────────────────────────────
\usepackage[a4paper, margin=2.5cm]{geometry}
\usepackage{graphicx}
\usepackage{booktabs}
\usepackage{longtable}
\usepackage{array}
\usepackage{float}
\usepackage{caption}
\usepackage{amsmath}
\usepackage{microtype}
\usepackage{parskip}
\usepackage{setspace}
\usepackage{hyperref}
\usepackage{xcolor}
\usepackage{enumitem}

\hypersetup{colorlinks=true, linkcolor=black, citecolor=black, urlcolor=black}
\setstretch{1.15}
\captionsetup{font=small, labelfont=bf, labelsep=period}

% ── document ──────────────────────────────────────────────────────────────────
\begin{document}

\begin{titlepage}
  \centering
  \vspace*{2cm}
  {\Large\bfseries Exploratory Data Analysis Report:\\[0.4em]
   Tropical Cyclone Dataset with Vedic-Astronomical Features}\\[1.5cm]
  {\large Comprehensive EDA Findings: Feature Importance, SHAP/LIME Explainability,\\
   Vedic Category Analysis, Retrograde Effects, and Circular Correlations}\\[1.0cm]
  {\normalsize Dataset: 47,533 observations across 944 tropical cyclones (2016--2026)\\
   186 features: 14 meteorological \textbullet\
   18 Vedic-Panchang \textbullet\ 154 astronomical}\\[2cm]
  {\small Generated by \texttt{09\_EDA\_report.py}}
\end{titlepage}

\tableofcontents
\newpage

% ══════════════════════════════════════════════════════════════════════════════
\section{Overview and Data Description}
% ══════════════════════════════════════════════════════════════════════════════

This report documents all major exploratory data analysis (EDA) findings from a
47,533-observation tropical cyclone dataset. Each observation represents one
6-hourly best-track position of a storm, enriched with three categories of features:

\begin{itemize}[noitemsep]
  \item \textbf{Meteorological (14 features):} latitude, longitude, wind speed, spline-derived
        track geometry, elevation, and land-proximity metrics.
  \item \textbf{Vedic-Panchang (18 features):} Tithi, Nakshatra, Yoga, Karana, Paksha,
        Moonsign, Samvatsara, day-length (Dinamana/Ratrimana), Weekday, Ayana, and Ritu.
  \item \textbf{Scientific-Astronomical (154 features):} Planet positions (ecliptic
        longitude/latitude, RA/Dec), altitudes, azimuths, angular separations (aspects),
        and Vedic degree representations with speed and padam.
\end{itemize}

\medskip
\noindent\textbf{Key variable:} Wind speed (kt) --- 10-minute maximum sustained surface wind.
Median = 40\,kt; mean $\approx$ 45\,kt; right-skewed (skewness $= 1.35$); all distributions
are non-normal (confirmed by Shapiro--Wilk and D'Agostino--Pearson tests), hence all
group comparisons use non-parametric methods throughout.

% ══════════════════════════════════════════════════════════════════════════════
\section{Feature Importance: Tree-Model Ensemble}
% ══════════════════════════════════════════════════════════════════════════════

Feature importances were computed using three tree-based models --- Random Forest,
XGBoost, and LightGBM --- trained with 5-fold GroupKFold cross-validation (groups
= \texttt{storm\_id}) to prevent within-storm autocorrelation leakage. Importances
were averaged across folds and models.

\medskip
\noindent\textbf{Key findings:}
\begin{itemize}[noitemsep]
  \item \texttt{longitude}: mean importance $= 0.0588$ (rank \#1)
  \item Top-5 features: \texttt{longitude}, \texttt{latitude}, \texttt{spline_end_lon}, \texttt{dinamana_minutes}, \texttt{vedic_chandra_speed_deg_per_day}.
  \item 10 Vedic features appear in the top-30 (out of 30 positions).
  \item Highest-ranked Vedic feature: \texttt{vedic_chandra_speed_deg_per_day} (mean importance $= 0.0203$).
  \item GroupKFold CV by \texttt{storm\_id} prevents within-storm leakage; negative Group CV $R^2$ on Vedic-only models confirms no out-of-storm generalisation.
\end{itemize}

\begin{figure}[H]
\centering
\includegraphics[width=0.92\linewidth]{C:/Users/reach/Downloads/drik_parser/old/v2/EDA/output/eda_report/feature_importance_top30.png}
\caption{Top-30 Feature Importances across RandomForest, XGBoost, and LightGBM. Geographic/track features (longitude, latitude, spline\_end\_lon) dominate. Vedic features appear from rank 5 onward.}
\label{fig:fi_top30}
\end{figure}


\begin{figure}[H]
\centering
\includegraphics[width=0.92\linewidth]{C:/Users/reach/Downloads/drik_parser/old/v2/EDA/output/eda_report/important_vs_unimportant.png}
\caption{Contrast between the 25 most important and 25 least important features (mean importance across RF/XGB/LGB). The bottom-25 features are overwhelmingly trigonometric altitude encodings and binary flags.}
\label{fig:fi_contrast}
\end{figure}



% ══════════════════════════════════════════════════════════════════════════════
\section{Mutual Information with Wind Speed}
% ══════════════════════════════════════════════════════════════════════════════

Mutual Information (MI) measures statistical dependence between each feature and
wind speed without assuming linearity. MI was computed using $k$-nearest-neighbour
estimation ($k=5$) on all 186 features.

\medskip
\noindent\textbf{Key findings:}
\begin{itemize}[noitemsep]
  \item Top feature: \texttt{lahiri_ayanamsha} (MI $= 0.720$).
  \item Top-5: \texttt{lahiri_ayanamsha}, \texttt{vedic_rahu_full_degree}, \texttt{vedic_rahu_right_ascension}, \texttt{vedic_ketu_right_ascension}, \texttt{vedic_ketu_full_degree}.
  \item Highest Vedic MI: \texttt{vedic_rahu_full_degree} (MI $= 0.691$) --- driven by temporal trend.
\end{itemize}

\noindent\textit{Note:} High MI for \texttt{lahiri\_ayanamsha} and Rahu/Ketu degree features
reflects their strong temporal trend (precession-corrected ecliptic longitude changes
systematically with year), which confounds any apparent predictive signal.

\begin{figure}[H]
\centering
\includegraphics[width=0.92\linewidth]{C:/Users/reach/Downloads/drik_parser/old/v2/EDA/output/eda_report/mutual_information_top30.png}
\caption{Top-30 features by Mutual Information with wind speed. Ayanamsha and Rahu/Ketu degree features lead due to temporal confounding. Meteorological features (latitude, spline geometry) follow.}
\label{fig:mi_top30}
\end{figure}



% ══════════════════════════════════════════════════════════════════════════════
\section{SHAP Explainability Analysis}
% ══════════════════════════════════════════════════════════════════════════════

SHAP (SHapley Additive exPlanations) values were computed using TreeExplainer on an
XGBoost model. Mean absolute SHAP values quantify each feature's average contribution
magnitude to individual wind speed predictions on the held-out test set.

\subsection{SHAP Feature Importance}

\begin{figure}[H]
\centering
\includegraphics[width=0.92\linewidth]{C:/Users/reach/Downloads/drik_parser/old/v2/EDA/output/eda_report/shap_summary_bar.png}
\caption{Mean absolute SHAP values for the top-30 features. Latitude dominates with mean $|\text{SHAP}| = 4.2$\,kt, followed by geographic features. Vedic-speed features (e.g., \texttt{vedic\_chandra\_speed}) appear in the top 10.}
\label{fig:shap_bar}
\end{figure}



\begin{figure}[H]
\centering
\includegraphics[width=0.92\linewidth]{C:/Users/reach/Downloads/drik_parser/old/v2/EDA/output/eda_report/shap_important_vs_unimportant.png}
\caption{Contrast: top-25 vs bottom-25 features by mean $|$SHAP$|$. Bottom features are predominantly trigonometric altitude encodings (e.g., \texttt{mars\_sin\_alt}, \texttt{saturn\_sin\_alt}) with near-zero SHAP values.}
\label{fig:shap_contrast}
\end{figure}



\subsection{SHAP Feature Interactions}

SHAP interaction values decompose each prediction into pairwise feature contributions.
The top interaction pairs are listed in Table~\ref{tab:shap_interact}.

\begin{table}[H]
  \centering
  \caption{Top-15 SHAP Feature Interactions by Mean Interaction Value}
  \label{tab:shap_interact}
  \begin{tabular}{lll}
    \toprule
    Feature 1 & Feature 2 & Mean Interaction \\
    \midrule
  latitude & spline_bearing_deg & 0.8220 \\
  latitude & spline_end_lon & 0.5601 \\
  longitude & spline_end_lon & 0.5514 \\
  longitude & latitude & 0.5339 \\
  latitude & ratrimana_minutes & 0.3631 \\
  latitude & vedic_varun_speed_deg_per_day & 0.2816 \\
  latitude & vedic_arun_speed_deg_per_day & 0.2572 \\
  latitude & sun_dec_deg & 0.2566 \\
  latitude & vedic_budha_full_degree & 0.2414 \\
  sun_dec_deg & mercury_dec_deg & 0.2374 \\
  latitude & vedic_guru_full_degree & 0.2361 \\
  vedic_budha_full_degree & dinamana_minutes & 0.2217 \\
  latitude & dinamana_minutes & 0.2174 \\
  latitude & vedic_chandra_right_ascension & 0.2067 \\
  sun_dec_deg & vedic_arun_speed_deg_per_day & 0.1856 \\
    \bottomrule
  \end{tabular}
\end{table}

\begin{figure}[H]
\centering
\includegraphics[width=0.92\linewidth]{C:/Users/reach/Downloads/drik_parser/old/v2/EDA/output/eda_report/shap_interactions_top20.png}
\caption{Top-20 SHAP feature interactions. Latitude $\times$ spline\_bearing\_deg is the dominant pair (interaction value $= 0.822$), reflecting how geographic position modulates storm track curvature effects on intensity.}
\label{fig:shap_interact}
\end{figure}



% ══════════════════════════════════════════════════════════════════════════════
\section{LIME Explainability and SHAP vs LIME Consensus}
% ══════════════════════════════════════════════════════════════════════════════

LIME (Local Interpretable Model-agnostic Explanations) was applied to five representative
instances spanning the wind speed distribution (10th--90th percentiles). Mean absolute
LIME weights were aggregated across instances to form a global LIME importance ranking.

SHAP and LIME importances were normalised to $[0, 1]$ and ranked jointly. Features
appearing in the top-30 of both methods are termed \textit{consensus features}.

\medskip
\noindent\textbf{Consensus findings:}
\begin{itemize}[noitemsep]
  \item 14 features appear in the top-30 of both SHAP and LIME.
  \item Consensus set: \textit{dinamana_minutes, latitude, longitude, mercury_dec_deg, nearest_lat, neptune_ecliptic_lon_deg, ratrimana_minutes, spline_bearing_deg, spline_end_lon, sun_dec_deg, vedic_arun_speed_deg_per_day, vedic_budha_full_degree}.
  \item Vedic features in consensus: vedic_budha_speed_deg_per_day, vedic_budha_full_degree, vedic_guru_full_degree, vedic_arun_speed_deg_per_day.
\end{itemize}

\begin{figure}[H]
\centering
\includegraphics[width=0.92\linewidth]{C:/Users/reach/Downloads/drik_parser/old/v2/EDA/output/eda_report/shap_vs_lime_top20.png}
\caption{Normalised SHAP vs LIME importance for the top-20 consensus features. Both methods agree on geographic and track features dominating. Disagreements appear for trajectory-geometry features (e.g., \texttt{spline\_curvature} ranks higher in SHAP than LIME).}
\label{fig:shap_vs_lime}
\end{figure}



% ══════════════════════════════════════════════════════════════════════════════
\section{Vedic Category Analysis: Kruskal--Wallis Tests}
% ══════════════════════════════════════════════════════════════════════════════

Non-parametric Kruskal--Wallis tests were applied to compare wind speed distributions
across all major Vedic categorical variables. Since all distributions are non-normal
(confirmed by Shapiro--Wilk), non-parametric methods are required.
Benjamini--Hochberg (BH) FDR correction was applied across all tests.

\begin{table}[H]
  \centering
  \caption{Kruskal--Wallis Test Results: Vedic Categories vs Wind Speed.
           BH-corrected significance threshold $\alpha = 0.05$.}
  \label{tab:kw}
  \begin{tabular}{lrrl r}
    \toprule
    Variable & $H$ Statistic & $p$-value & \# Categories & Significant \\
    \midrule
  drik_ritu & 1214.03 & 2.68e-260 & 6 & Yes \\
  sunsign & 1123.14 & 5.89e-234 & 12 & Yes \\
  yoga & 143.35 & 3.43e-18 & 27 & Yes \\
  nakshatra & 116.09 & 2.37e-13 & 27 & Yes \\
  drik_ayana & 102.63 & 4.04e-24 & 2 & Yes \\
  moonsign & 77.93 & 3.71e-12 & 12 & Yes \\
  tithi & 51.26 & 7.50e-06 & 16 & Yes \\
  karana & 42.48 & 6.15e-06 & 11 & Yes \\
  weekday & 11.59 & 0.0717 & 7 & No \\
  paksha & 8.11 & 0.0044 & 2 & Yes \\
    \bottomrule
  \end{tabular}
\end{table}

\medskip
\noindent\textbf{Critical caveats:}
\begin{itemize}[noitemsep]
  \item \texttt{drik\_ritu} ($H=1214$) and \texttt{sunsign} ($H=1123$) are \emph{calendar-season
        proxies}, not independent Vedic signals. Their significance reflects the well-known
        seasonal cycle of tropical cyclone activity.
  \item All tests use $N = 47,533$ observations. With such large $N$,
        even tiny group differences ($<0.5$\,kt) achieve $p \ll 0.001$. Effect sizes must
        be consulted before any practical interpretation.
  \item Confound-adjusted analyses (partial KW within basin $\times$ season strata) show
        that Yoga, Nakshatra, Tithi, and Moonsign retain significance in 100\% of strata ---
        however, the effect sizes ($\varepsilon^2 < 0.003$) remain negligible.
\end{itemize}

% ══════════════════════════════════════════════════════════════════════════════
\section{Nakshatra Intensity Analysis}
% ══════════════════════════════════════════════════════════════════════════════

Median wind speed was computed for each of the 27 Nakshatras across all observations.
Differences between Nakshatras are statistically significant ($H = 116$, $p = 2.4 \times 10^{-13}$)
but the effect size (epsilon-squared $\varepsilon^2 \approx 0.002$) is trivially small.

\begin{figure}[H]
\centering
\includegraphics[width=0.92\linewidth]{C:/Users/reach/Downloads/drik_parser/old/v2/EDA/output/eda_report/nakshatra_intensity_bar.png}
\caption{Median wind speed (kt) per Nakshatra, sorted descending with IQR error bars. Numbers above bars indicate observation count per Nakshatra. The $\approx 4$\,kt range across Nakshatras is statistically significant but practically negligible ($\varepsilon^2 \approx 0.002$).}
\label{fig:nakshatra}
\end{figure}



% ══════════════════════════════════════════════════════════════════════════════
\section{Tithi $\times$ Nakshatra Interaction Heatmaps}
% ══════════════════════════════════════════════════════════════════════════════

Cross-tabulation of Tithi (lunar day, 1--30) and Nakshatra (27 lunar mansions)
reveals the joint distribution of mean and median wind speeds across these two
fundamental Vedic time categories.

\begin{figure}[H]
\centering
\includegraphics[width=0.92\linewidth]{C:/Users/reach/Downloads/drik_parser/old/v2/EDA/output/eda_report/tithi_nakshatra_heatmap.png}
\caption{Mean wind speed (kt) by Tithi (rows) and Nakshatra (columns). Darker cells indicate higher mean wind. Cells with no observations appear white. The pattern reflects the co-occurrence structure of the Vedic calendar rather than independent effects.}
\label{fig:tn_mean}
\end{figure}



\begin{figure}[H]
\centering
\includegraphics[width=0.92\linewidth]{C:/Users/reach/Downloads/drik_parser/old/v2/EDA/output/eda_report/tithi_nakshatra_heatmap_median.png}
\caption{Median wind speed (kt) by Tithi (rows) and Nakshatra (columns). More robust to outliers than the mean heatmap. No systematic high-intensity zone is visible across any Tithi--Nakshatra combination.}
\label{fig:tn_median}
\end{figure}



% ══════════════════════════════════════════════════════════════════════════════
\section{Retrograde Planetary Motion and Storm Intensity}
% ══════════════════════════════════════════════════════════════════════════════

For each of the seven classical Vedic planets (Guru/Jupiter, Shani/Saturn,
Mangal/Mars, Budha/Mercury, Shukra/Venus, Arun/Uranus, Varun/Neptune),
observations were split into retrograde (negative geocentric speed) and
direct (positive geocentric speed) periods. Mann--Whitney U tests compared
wind speed distributions.

\begin{table}[H]
  \centering
  \caption{Retrograde vs Direct Planetary Motion: Mann--Whitney U Results.
           $N_{\text{retro}}$ and $N_{\text{direct}}$ sum to $\approx$ 47,533.
           Wind difference = median(retrograde) $-$ median(direct).
           ${}^*$ = significant at $p < 0.05$.}
  \label{tab:retro}
  \begin{tabular}{lrrlrrrr}
    \toprule
    Planet & $N_\text{retro}$ & Retro\% & $p$ & Med (retro) & Med (direct) & $\Delta$Med & Sig. \\
    \midrule
  guru & 14,128 & 29.7\% & 40 & 42 & -2.0 & 2.31e-23 & \textbf{*} \\
  shani & 22,023 & 46.3\% & 40 & 40 & +0.0 & 4.99e-06 & \textbf{*} \\
  mangal & 4,317 & 9.1\% & 41 & 40 & +1.0 & 0.0004 & \textbf{*} \\
  budha & 9,889 & 20.8\% & 40 & 40 & +0.0 & 2.34e-22 & \textbf{*} \\
  shukra & 3,749 & 7.9\% & 45 & 40 & +5.0 & 6.62e-06 & \textbf{*} \\
  arun & 24,621 & 51.8\% & 43 & 40 & +3.0 & 1.02e-38 & \textbf{*} \\
  varun & 28,665 & 60.3\% & 43 & 39 & +4.0 & 2.33e-59 & \textbf{*} \\
    \bottomrule
  \end{tabular}
\end{table}

\medskip
\noindent\textbf{Interpretation:}
All seven planets show statistically significant differences in wind speed
between retrograde and direct periods ($p < 0.001$). However, the magnitude
of differences is small ($\Delta\text{median} \leq 4$\,kt in all cases).
The effect sizes (rank-biserial $r < 0.10$ for all planets) indicate negligible
practical significance. With $N > 28{,}000$ per group, even 1\,kt differences
achieve $p \ll 0.001$.

\begin{figure}[H]
\centering
\includegraphics[width=0.92\linewidth]{C:/Users/reach/Downloads/drik_parser/old/v2/EDA/output/eda_report/retrograde_effects.png}
\caption{Retrograde vs direct median wind speed per planet. Asterisks ($^*$) indicate statistical significance ($p < 0.05$). All differences are $\leq 4$\,kt. Arun and Varun show the largest differences, but these outer planets are retrograde for $>50\%$ of observations, making the groups unequal.}
\label{fig:retrograde}
\end{figure}



% ══════════════════════════════════════════════════════════════════════════════
\section{Planetary Aspect Correlations}
% ══════════════════════════════════════════════════════════════════════════════

Planetary aspects are angular separations between planet pairs (modulo 360\textdegree).
Pearson correlations between each aspect and wind speed were computed across all
47,533 observations.

\begin{table}[H]
  \centering
  \caption{Planetary Aspect Pearson Correlations with Wind Speed (all 47,533 observations).
           Maximum $|r| = 0.0766$.}
  \label{tab:aspects}
  \begin{tabular}{lr}
    \toprule
    Aspect Pair & Pearson $r$ \\
    \midrule
  aspect_jupiter_saturn & 0.076636 \\
  aspect_moon_mars & 0.033083 \\
  aspect_sun_saturn & 0.028431 \\
  aspect_mercury_saturn & 0.023740 \\
  aspect_moon_jupiter & 0.015221 \\
  aspect_sun_mercury & 0.014608 \\
  aspect_venus_saturn & 0.013318 \\
  aspect_mercury_mars & 0.007700 \\
  aspect_mars_saturn & 0.004658 \\
  aspect_moon_saturn & 0.004187 \\
  aspect_sun_mars & 0.001534 \\
  aspect_moon_mercury & -0.000170 \\
  aspect_sun_moon & -0.000202 \\
  aspect_moon_venus & -0.002299 \\
  aspect_mercury_venus & -0.007633 \\
  aspect_mars_jupiter & -0.009821 \\
  aspect_sun_venus & -0.010910 \\
  aspect_mercury_jupiter & -0.016471 \\
  aspect_sun_jupiter & -0.017351 \\
  aspect_venus_mars & -0.025494 \\
  aspect_venus_jupiter & -0.039407 \\
    \bottomrule
  \end{tabular}
\end{table}

\begin{figure}[H]
\centering
\includegraphics[width=0.92\linewidth]{C:/Users/reach/Downloads/drik_parser/old/v2/EDA/output/eda_report/aspect_correlations.png}
\caption{Pearson correlations of planetary aspect angles with wind speed. The Jupiter--Saturn aspect shows the highest positive correlation ($r = 0.077$), and Venus--Jupiter the most negative ($r = -0.039$). All correlations are weak ($|r| < 0.10$).}
\label{fig:aspects}
\end{figure}



\medskip
\noindent\textbf{Interpretation:} The maximum correlation is $r \approx 0.077$ (Jupiter--Saturn
separation), explaining less than $0.6\%$ of variance in wind speed. These are the
weakest signals in the entire dataset.

% ══════════════════════════════════════════════════════════════════════════════
\section{Vedic Planetary Degree: Circular-Linear Correlations}
% ══════════════════════════════════════════════════════════════════════════════

Vedic planet degrees are cyclical variables (0--360\textdegree). Naive Pearson
correlation on raw degrees is \emph{incorrect} because it treats 1\textdegree\ and
359\textdegree\ as far apart. The correct method decomposes the degree angle into
$\sin(\theta)$ and $\cos(\theta)$ components and reports
$r_{\text{circ}} = \max(|r_{\sin}|, |r_{\cos}|)$.

\begin{table}[H]
  \centering
  \caption{Circular-Linear Correlations: Vedic Planet Degrees vs Wind Speed.
           $r_{\sin}$ = Pearson correlation of $\sin(\text{degree})$ with wind;
           $r_{\cos}$ = Pearson correlation of $\cos(\text{degree})$ with wind;
           $r_{\text{circ}}$ = $\max(|r_{\sin}|, |r_{\cos}|)$.}
  \label{tab:vedic_circ}
  \begin{tabular}{lrrr}
    \toprule
    Feature & $r_{\sin}$ & $r_{\cos}$ & $r_{\text{circ}}$ \\
    \midrule
  vedic_yam_full_degree & 0.08938 & -0.02018 & 0.08938 \\
  vedic_guru_full_degree & 0.07130 & -0.04249 & 0.07130 \\
  vedic_shani_full_degree & 0.05394 & -0.01491 & 0.05394 \\
  vedic_spashth_ketu_full_degree & 0.02863 & 0.03457 & 0.03457 \\
  vedic_ketu_full_degree & 0.02739 & 0.03392 & 0.03392 \\
  vedic_lagna_full_degree & -0.00098 & 0.00372 & 0.00372 \\
  vedic_chandra_full_degree & -0.00255 & -0.00875 & -0.00875 \\
  vedic_varun_full_degree & -0.00258 & -0.00899 & -0.00899 \\
  vedic_arun_full_degree & 0.00522 & -0.01990 & -0.01990 \\
  vedic_rahu_full_degree & -0.02739 & -0.03392 & -0.03392 \\
  vedic_spashth_rahu_full_degree & -0.02863 & -0.03457 & -0.03457 \\
  vedic_mangal_full_degree & -0.00326 & -0.05290 & -0.05290 \\
  vedic_budha_full_degree & 0.02773 & -0.05700 & -0.05700 \\
  vedic_surya_full_degree & 0.04194 & -0.06041 & -0.06041 \\
  vedic_shukra_full_degree & 0.02630 & -0.06712 & -0.06712 \\
    \bottomrule
  \end{tabular}
\end{table}

\begin{figure}[H]
\centering
\includegraphics[width=0.92\linewidth]{C:/Users/reach/Downloads/drik_parser/old/v2/EDA/output/eda_report/vedic_degree_corr_bar.png}
\caption{Circular-linear correlations $r_{\text{circ}}$ per Vedic planet degree. Yama (\texttt{vedic\_yam}) shows the highest correlation ($r \approx 0.089$), followed by Guru/Jupiter ($r \approx 0.071$). All correlations are below $0.10$ --- explaining less than $1\%$ of wind speed variance.}
\label{fig:vedic_circ}
\end{figure}



% ══════════════════════════════════════════════════════════════════════════════
\section{Variance Inflation Factor (VIF) Analysis}
% ══════════════════════════════════════════════════════════════════════════════

VIF analysis was applied to a representative subset of 14 features to check for
multicollinearity. VIF values $> 5$ indicate moderate, and $> 10$ indicate severe
multicollinearity.

\begin{table}[H]
  \centering
  \caption{Variance Inflation Factors for Representative Features.
           All values $< 2.0$ indicate low multicollinearity in this subset.}
  \label{tab:vif}
  \begin{tabular}{lrl}
    \toprule
    Feature & VIF & Assessment \\
    \midrule
  latitude & 1.992 & Low \\
  saturn_altitude_deg & 1.709 & Low \\
  day_of_year & 1.703 & Low \\
  jupiter_altitude_deg & 1.579 & Low \\
  sun_ecliptic_lon_deg & 1.481 & Low \\
  dinamana_minutes & 1.471 & Low \\
  longitude & 1.214 & Low \\
  sun_altitude_deg & 1.096 & Low \\
  hour_utc & 1.087 & Low \\
  elevation & 1.079 & Low \\
  sun_azimuth_deg & 1.029 & Low \\
  moon_ecliptic_lon_deg & 1.010 & Low \\
  moon_altitude_deg & 1.008 & Low \\
  moon_azimuth_deg & 1.006 & Low \\
    \bottomrule
  \end{tabular}
\end{table}

\noindent\textbf{Finding:} All VIF values are below 2.0, indicating negligible multicollinearity in this feature subset. 
Note that this subset was selected to represent diverse feature types; the full
186-feature space contains groups of near-duplicate features (e.g., Rahu vs Spashth-Rahu)
with higher inter-group correlations.

% ══════════════════════════════════════════════════════════════════════════════
\section{Intensity Change Analysis}
% ══════════════════════════════════════════════════════════════════════════════

Wind change per 6-hour step ($\Delta\text{wind}$) was computed from consecutive
best-track positions within each storm. Rapid Intensification (RI) is defined as
$\Delta\text{wind} \geq +30$\,kt per 6 hours; Rapid Weakening (RW) as
$\Delta\text{wind} \leq -30$\,kt per 6 hours.

\begin{figure}[H]
\centering
\includegraphics[width=0.92\linewidth]{C:/Users/reach/Downloads/drik_parser/old/v2/EDA/output/eda_report/intensity_change.png}
\caption{Intensity change analysis. Left: distribution of 6-hourly wind changes with RI ($\geq +30$\,kt) and RW ($\leq -30$\,kt) thresholds marked. Centre: counts of RI and RW events vs total observations. Right: scatter of wind change vs sun altitude angle (random 5000-sample subset).}
\label{fig:intensity}
\end{figure}



% ══════════════════════════════════════════════════════════════════════════════
\section{Summary and Methodological Notes}
% ══════════════════════════════════════════════════════════════════════════════

\subsection{Feature Importance Hierarchy}

Across all importance metrics (RF/XGB/LGB mean importance, SHAP, LIME, Mutual
Information), the feature hierarchy is consistent:

\begin{enumerate}[noitemsep]
  \item \textbf{Geographic position:} Latitude ($\#1$ in SHAP, LIME, and RF importance),
        longitude, and spline endpoint coordinates.
  \item \textbf{Track geometry:} Spline bearing, curvature, displacement.
  \item \textbf{Day-length timing:} \texttt{dinamana\_minutes} and \texttt{ratrimana\_minutes}
        (seasonal/latitude proxies).
  \item \textbf{Planet speeds and degrees:} Moon speed, Mercury degree/speed, Guru/Jupiter
        degree appear in the top 20 across multiple methods.
  \item \textbf{Altitude/azimuth features:} Consistently unimportant across all methods.
\end{enumerate}

\subsection{Vedic Feature Signals}

No Vedic-Panchang feature appears in the top-5 by any importance metric. In the top-30:
10 Vedic features appear (out of 30 positions)
using mean tree-model importance. Vedic features that do appear are predominantly
\emph{speed} features (Moon speed, Mercury speed, Uranus speed) which co-vary with
geographic position and season.

\subsection{Large-$N$ Significance Inflation}

All statistical tests were conducted on $N = 47,533$ observations.
At this sample size, even differences of $< 1$\,kt achieve $p < 0.001$.
Effect sizes must be consulted for all significant $p$-values:

\begin{itemize}[noitemsep]
  \item Kruskal--Wallis: $\varepsilon^2 < 0.003$ for all Vedic categorical variables.
  \item Retrograde effects: rank-biserial $r < 0.10$ for all planets.
  \item Circular correlations: $r_{\text{circ}} < 0.09$ for all Vedic degrees.
  \item Aspect correlations: $|r| < 0.08$ for all planetary pairs.
\end{itemize}

\subsection{Autocorrelation}

Consecutive observations within the same storm are temporally autocorrelated
(ACF at lag-1 $\approx 0.92$ for wind speed). All group-level tests treat each
observation independently. Group cross-validation (GroupKFold by storm\_id) was
used for all predictive models to prevent leakage, but statistical tests
(KW, Mann--Whitney) remain subject to this limitation.

% ══════════════════════════════════════════════════════════════════════════════
\section{Conclusion}
% ══════════════════════════════════════════════════════════════════════════════

The EDA confirms that:
\begin{enumerate}[noitemsep]
  \item \textbf{Geographic features dominate} all importance metrics, consistent with
        known climatological structure of tropical cyclone intensity.
  \item \textbf{Vedic-Panchang categorical variables} show statistically significant but
        practically negligible wind speed differences across their categories. The
        dominant confound is the seasonal cycle (Ritu, Ayana).
  \item \textbf{Retrograde planetary motion} is associated with statistically significant
        but very small ($\leq 4$\,kt) median wind speed differences across all 7 planets.
  \item \textbf{Circular-linear correlations} of Vedic planet degrees with wind speed are
        below $r = 0.09$ for all planets --- explaining $< 1\%$ of variance.
  \item \textbf{Planetary aspects} have negligible correlations with wind speed ($|r| < 0.08$).
  \item \textbf{Multicollinearity} among the 14 analysed features is low (all VIF $< 2$),
        though feature groups (Rahu/Ketu variants, speed/degree pairs) exhibit redundancy.
  \item \textbf{SHAP--LIME consensus} on 14 features
        provides robust interpretability; both methods agree that Vedic features are marginal
        contributors to model predictions.
\end{enumerate}

\end{document}
